\documentclass[twocolumn]{article}

\usepackage[T1]{fontenc}

\DeclareRobustCommand{\VAN}[3]{#2}
\let\VANthebibliography\thebibliography
\def\thebibliography{\DeclareRobustCommand{\VAN}[3]{##3}\VANthebibliography}


\usepackage{graphicx}	% Including figure files
\usepackage{amsmath}	% Advanced maths commands
\usepackage{amssymb}	% Extra maths symbols

\usepackage{tabularx}
\usepackage{comment}
\usepackage{float}
\usepackage{xcolor}

\usepackage[top=3cm,bottom=3cm,left=2cm,right=2cm]{geometry}

\setlength{\columnsep}{3em}

\newcommand{\han}{\textcolor{orange}}
\newcommand{\fref}{f_{\rm ref}}
\newcommand{\ARed}{A_{\rm sn}}
\newcommand{\gammaRed}{\gamma_{\rm sn}}
\newcommand{\ARedI}{A_{{\rm sn},I}}
\newcommand{\ARedJ}{A_{{\rm sn},J}}
\newcommand{\gammaRedI}{\gamma_{{\rm sn},I}}
\newcommand{\gammaRedJ}{\gamma_{{\rm sn},J}}
\newcommand{\Pg}{P_{\rm g}}
\newcommand{\PI}{P_I}
\newcommand{\PJ}{P_J}
\newcommand{\PsnI}{P_{{\rm sn},I}}
\newcommand{\chiIJhd}{\chi_{IJ}}
\newcommand{\chiIJdip}{\chi_{IJ}^{\rm dip}}
\newcommand{\sigJitter}{\sigma_{\rm jit}}

\input{abbreviations.tex}

\title{MeerTime pulsar timing array observation strategy}

\begin{document}
\maketitle

MeerTime is the pulsar timing program of the MeerKAT radio observatory. 
One of the MeerTime projects is a pulsar timing array (PTA) campaign to search for nanohertz gravitational waves from massive black hole binaries using around $10$ years of data. 
This project is focusing on improving the PTA's signal to noise ratio (SNR) for the stochastic gravitational wave background emitted by a population of massive black hole binaries. 



\section{Average signal to noise ratio}

The average S/N for a given PTA and a predicted  SGWB signal following the method in~\cite{SiemensEtAl:2013}. 
The power spectrum for a SGWB is given by
\begin{equation}
\Pg (f) = \frac{A^2}{24 \pi^2} \left(\frac{f}{\fref}\right)^{2\alpha} f^{-3}\,,
\label{eqn:Pgw}
% equation 16 in SiemansEtAl:2013
\end{equation}
where $f$ is the gravitational-wave frequency, $\fref$ is a reference frequency chosen as $1\,{\rm yr}^{-1}$, $A$ is the amplitude of the background at $f=\fref$, and $\alpha$ is the slope of the spectrum. 
Throughout this work we assume $a=-2/3$ consistent with a stochastic background from circular massive black hole binaries and $A = 2 \times 10^{-15}$ consistent with recent speculative hints of detection~\citep{GoncharovEtAlPPTA:2021,NANOGrav12.5GWB:2020}.

The average SNR for a PTA with total timespan $T$ is given by
\begin{equation}
\left< \rho \right> = \left( 2 T \sum_{IJ} \chiIJhd^2 
                             \int_{f_L}^{f_H} {\rm d}f  
                             \frac{\Pg^2 (f)}{\PI(f)\PJ(f)}
                            \right)^{1/2} \,,
\label{eqn:snrPIPJPg}
\end{equation}
where $\sum_{IJ}$ is the sum over all pulsar pairs ($I \neq J$) and $P_I$ is the power spectrum of the $I{\rm th}$ pulsar. 
We assume that the timespan of observations, $T$ is the same for all pulsars. 
The integration limits are given by $f_L=1/T$ and $f_H=c/2$ where $c$ is the cadence of observations, which we also assume to be the identical for all pulsars in the array. 
The Hellings and Downs coefficients $\chi_{IJ}$ describe the expected angular correlation produced by the SGWB signal in pairs of pulsars and is given by~\citep{HellingsDowns:1983}
\begin{eqnarray}
\chiIJhd &=& \frac{1}{2} 
           - \frac{1}{4} \left(\frac{1 - \cos(\zeta)}{2}\right) \nonumber\\
           & &+ \frac{3}{2} \left(\frac{1 - \cos(\zeta)}{2}\right) 
           \log \left( \frac{1 - \cos(\zeta)}{2} \right)\,,
\label{eqn:hd}
\end{eqnarray}
for pulsars separated by angle $\zeta$.

If we assume evenly sampled data and that there is only white noise in the pulsars then $\PI$ is the sum of Eqn.~\ref{eqn:Pgw} plus white noise, 
\begin{equation}
\PI (f) = \frac{A^2}{24 \pi^2} \left(\frac{f}{\fref}\right)^{2\alpha} f^{-3} 
          + 2 \sigma_I^2 \Delta t\,,
\end{equation}
where $c=1/\Delta t$ and $\sigma_I$ is the timing precision for the $I{\rm th}$ pulsar. 

Pulsars may also contain spin noise (or red noise) and jitter noise due to changes in the pulse profile shape. 
The power spectral density of spin noise for the $I{\rm th}$ pulsar is given by 
\begin{equation}
% from Goncharov 2021 eqn 5
\PsnI (f) = \frac{\ARedI^2}{12 \pi^2} 
            \left( \frac{f}{\fref} \right)^{-\gammaRedI} {\rm yr}^3\,,
\end{equation}
where $\ARedI$ and $\gammaRedI$ are the amplitude and slope of the power law for the $I$th pulsar and ${\rm yr}$ is the number of seconds in a year. 
Jitter noise $\sigma_{\rm jit}$ is added in quadrature to the white noise part the the spectrum.
The resulting power spectrum is
\begin{eqnarray}
\PI (f) &=& \frac{A^2}{24 \pi^2} \left(\frac{f}{\fref}\right)^{2\alpha} f^{-3} \nonumber\\
          & &~+~ 2 \left( \sigma_I^2 + \sigma_{{\rm jit}, I}^2 \right) \Delta t \nonumber\\
          & &~+~ \frac{\ARedI^2}{12 \pi^2} 
          \left( \frac{f}{\fref} \right)^{-\gammaRedI} {\rm yr}^3\,,
\label{eqn:PIGWWR}
\end{eqnarray}
which includes SGWB, white noise, spin noise, and jitter noise.
Equations~\ref{eqn:snrPIPJPg} and~\ref{eqn:PIGWWR} can therefore be used to compute the average PTA SNR for a given observation strategy and assumed SGWB signal. 
\begin{comment}
\begin{equation}
\left< \rho \right> = \left( 2 T \sum_{IJ} \chi_{IJ}^2
                             \int_{f_L}^{f_H} {\rm d}f  
                             \frac{\Pg^2 (f)}
                                  {[\Pg (f) + 2\sigma_I^2\Delta t + \Psn (f)]
                                   [\Pg (f) + 2\sigma_J^2\Delta t + \Psn (f)]} 
                            \right)^{1/2} \,.
\label{eqn:snrWhiteNoiseOnly}
\end{equation}
\end{comment}




\section{Data}

MeerTime is timing $\NMeerTimePSRs$ pulsars.
Using observations to date, an anticipated $\sigma_I$ and integration time (for a single pulsar observation) required to achieve it can be estimated for each pulsar.
The value of $\sigma_I$ scales with observation $T_{{\rm obs},I}$ time like 
\begin{equation}
\sigma_I = \delta_I / \sqrt{T_{{\rm obs},I}}\,,
\label{eqn:TobsToSig}
\end{equation}
where $\delta_I$ is a constant of proportionality which depends on the pulsar. 
Throughout this work we make the simplifying assumption that $\delta_I$ is a constant for each pulsar throughout the PTA lifespan \han{[maybe add some discussion here about this?]}

\subsection{Spin noise and jitter noise}
MeerTime observations are not currently long enough to estimate spin noise in every pulsar.
However spin noise measurements already exist for $13$ of the pulsars from PPTA and NANOGrav analysis. 
We take $\ARed$ and $\gammaRed$ values from ~\cite{GoncharovEtAl:2021} (sixth and seventh column in Table 1) and ~\cite{NANOGrav12p5yr:2021} (tenth and eleventh column in Table 5). 
Where a value is listed in both papers, we take the~\cite{GoncharovEtAl:2021} value. 
Table~\ref{tab:redNoise} shows the red noise parameters for $13$ of the $\NMeerTimePSRs$ pulsars. 
Otherwise, we assume zero spin noise. 

\begin{table}
\begin{center}
\begin{tabular}{llll}
\hline \hline
Pulsar     & $\log_{10}(\ARed)$  & $\gammaRed$ & Reference \\
\hline
J0030+0451 & $-15.72$ & $5.3$   &  \cite{NANOGrav12p5yr:2021}\\
J0437-4715 & $-14.56$ & $2.99$  &  \cite{GoncharovEtAl:2021}\\
J0613-0200 & $-14.26$ & $4.17$  &  \cite{GoncharovEtAl:2021}\\
J0711-6830 & $-13.04$ & $1.09$  &  \cite{GoncharovEtAl:2021}\\
J1024-0719 & $-14.62$ & $6.39$  &  \cite{GoncharovEtAl:2021}\\
J1600-3053 & $-14.34$ & $3.81$  &  \cite{GoncharovEtAl:2021}\\
J1643-1224 & $-12.85$ & $0.98$  &  \cite{GoncharovEtAl:2021}\\
J1713+0747 & $-15.16$ & $1.6$   &  \cite{NANOGrav12p5yr:2021}\\
J1744-1134 & $-14.16$ & $1.7$   &  \cite{NANOGrav12p5yr:2021}\\
J1747-4036 & $-13.78$ & $3.8$   &  \cite{NANOGrav12p5yr:2021}\\
J1909-3744 & $-14.74$ & $4.05$  &  \cite{GoncharovEtAl:2021}\\
J2145-0750 & $-13.86$ & $1.6$   &  \cite{NANOGrav12p5yr:2021}\\
J2317+1439 & $-16.50$ & $6.0$   &  \cite{NANOGrav12p5yr:2021}\\
\hline\hline
\end{tabular}
\end{center}
\caption{\label{tab:redNoise}
Spin noise parameters for $13$ pulsars. 
The second and third columns show the log of the spin noise amplitude and the slope. 
}
\end{table}

Some of the pulsars have jitter noise estimates. 
Table~\ref{tab:jitter} shows jitter noise measurements for $13$ MeerTime pulsars from~\cite{ParthasarathyEtAl:2021}.
We assume the rest of the pulars have zero jitter noise.


\begin{table}
\begin{center}
\begin{tabular}{lll}
\hline \hline
Pulsar     & $\sigJitter$   & Reference \\
\hline
J0125-2327 & $48  \pm 13$   & ~\cite{ParthasarathyEtAl:2021}\\
J0437-4715 & $50  \pm 10$   & ~\cite{ParthasarathyEtAl:2021}\\
J0636-3044 & $100 \pm 30$   & ~\cite{ParthasarathyEtAl:2021}\\
J0711-6830 & $60  \pm 20$   & ~\cite{ParthasarathyEtAl:2021}\\
J1022+1001 & $120 \pm 20$   & ~\cite{ParthasarathyEtAl:2021}\\
J1045-4509 & $130 \pm 75$   & ~\cite{ParthasarathyEtAl:2021}\\
J1603-7202 & $180 \pm 40$   & ~\cite{ParthasarathyEtAl:2021}\\
J1730-2304 & $80  \pm 45$   & ~\cite{ParthasarathyEtAl:2021}\\
J1744-1134 & $30  \pm  6$   & ~\cite{ParthasarathyEtAl:2021}\\
J1757-5322 & $130 \pm 45$   & ~\cite{ParthasarathyEtAl:2021}\\
J1909-3744 & $9   \pm  3$   & ~\cite{ParthasarathyEtAl:2021}\\
J2145-0750 & $200 \pm 20$   & ~\cite{ParthasarathyEtAl:2021}\\
J2241-5236 & $3.8 \pm  0.8$ & ~\cite{ParthasarathyEtAl:2021}\\
\hline\hline
\end{tabular}
\end{center}
\caption{\label{tab:jitter}
Jitter noise measurements for $13$ pulsars. 
}
\end{table}

% pulsar reference - jitter, white nooise, red noies, position

\begin{comment}
% value from NANOGRav table is An
% An  = 0.006 \mu s yr^{1/2} 
% An -> 0.006 * 3.17E-14 yr^{3/2}
\begin{table}
\begin{tabular}{llll}
\hline\hline
Pulsar     & $\ARed$ & $\gammaRed$ & Ref.\\
\hline
J0030+0451 & $-15.72$ & $5.3$   &  \cite{NANOGrav12p5yr:2021}\\
J0437-4715 & $-14.56$ & $2.99$  &  \cite{GoncharovEtAl:2021}\\
J0613-0200 & $-14.26$ & $4.17$  &  \cite{GoncharovEtAl:2021}\\
J0711-6830 & $-13.04$ & $1.09$  &  \cite{GoncharovEtAl:2021}\\
J1024-0719 & $-14.62$ & $6.39$  &  \cite{GoncharovEtAl:2021}\\
J1600-3053 & $-14.34$ & $3.81$  &  \cite{GoncharovEtAl:2021}\\
J1643-1224 & $-12.85$ & $0.98$  &  \cite{GoncharovEtAl:2021}\\
J1713+0747 & $-15.16$ & $1.6$   &  \cite{NANOGrav12p5yr:2021}\\
J1744-1134 & $-14.16$ & $1.7$   &  \cite{NANOGrav12p5yr:2021}\\
J1747-4036 & $-13.78$ & $3.8$   &  \cite{NANOGrav12p5yr:2021}\\
J1909-3744 & $-14.74$ & $4.05$  &  \cite{GoncharovEtAl:2021}\\
J2145-0750 & $-13.86$ & $1.6$   &  \cite{NANOGrav12p5yr:2021}\\
J2317+1439 & $-16.50$ & $6.0$   &  \cite{NANOGrav12p5yr:2021}\\
\hline
\end{tabular}
\caption{\label{tab:redNoise}
Spin noise parameters for $13$ pulsars. 
[Notes: Values from Ref.~\cite{GoncharovEtAl:2021} are taken directly form their Table 1. 
%Values from Ref.~\cite{NANOGrav12p5yr:2021} are converted from units of ${\rm \mu s \,yr}^{1/2}$ to $\log_{10} \ARed[{\rm yr}^{3/2}]$.]
}
\end{table}
\end{comment}

\begin{comment}
\textbf{estimating spin noise for other pulsars? }
Spin noise parameters for the remaining pulsars can be estimated using the scaling relation formalism in Shannon and Cordes. 
The spin noise RMS value is given by 
\begin{equation}
\sigma_{\rm SN} = C_2 \mu^{\alpha} | {\dot \mu}_{-15} |^{\beta} t_{\rm yr}^{\gamma}
\end{equation}
where $\mu$......
In this work we use the central fit values for the ${\rm MSP}_{\rm 10,PPTA}$+${\rm NANO}$ ($\{ \ln C_2, \alpha, \beta, \gamma \} = \{-1.2,-0.9,0.8,2.3\}$).
\end{comment}




\subsection{Current observation strategy}
MeerTime pulsars are observed with a cadence of once every two weeks ($c=26\,{\rm yr}^{-1})$. 
The total observation time used in each $2$ week observing time is $\TwoWkObsTimeS\,{\rm s} = \TwoWkObsTimeHrs\,{\rm hrs}$, which is divided between the MeerTime pulsars. 
The targeted precision for each pulsar is $\sigma_I = 1 {\rm \mu\, s}$. 

Using the current observation strategy and assuming a GWB signal with $A = 2\times 10^{-15}$ and $\alpha = -2/3$, we can estimate the predicted average SNR for MeerTime from Eqn.~\ref{eqn:snrPIPJPg}. 
After a timespan of $T=10\,{\rm yrs}$, the average SNR for the current obseravation stategy is $12.3$.
Here, we consider whether, or not, the average SNR can be improved by a re-allocation of observing time between pulsars. 






\section{Initial attempts}


In this work we take a greedy algorithm approach to search for SNR improvements. 
The approach has its limitations in that it will always take the locally optimal step, and can miss the global optimal solution. 
However, it has a reasonably low computational cost and can still achieve SNR improvement. 


The method we use focuses on reorganising the time allocation between pulsars. 
We use the initial integration times as a starting point (see Section~\ref{sec:data}).
The total observation time available to MeerTime is constant ($\sum_{I} T_{{\rm obs},I} \TwoWkObsTimeS\,{\rm s}$) and the cadence is fixed at $c=26\,{\rm yr}^{-1}$.
Integration time is shuffled between pulsars to find an SNR improvement. 
At each step, a fraction $f$ of integration time is subtracted from a pulsar and reallocated to another pulsar so that 
\begin{eqnarray}
T_{{\rm obs},I} &\rightarrow& T_{{\rm obs},I} - f T_{{\rm obs},I}\,, \nonumber\\
T_{{\rm obs},J} &\rightarrow& T_{{\rm obs},J} + f T_{{\rm obs},I}\,, \nonmumber
\end{eqnarray}
where we choose $f=1/8$. 
Every possible combination of time shuffles between pulsars is trialled.
The combination which gives the biggest SNR increase at $10\,{\rm yrs}$ is chosen, providing an updated list of $T_{{\rm obs},I}$.
The new time allocations become the starting point for the next time shuffle. 
The process is repeated until the algorithm cannot find further gains in SNR. 

The shape of the Hellings-Downs function (equation~\ref{}) means that pulsars pairs with small angular separations have higher contributions to the SNR in equation~\ref{eqn:snrPIPJPg}. 
This leads to pulsars in the galactic centre being favoured in the time shuffle due to a higher number of small angular separations 

\subsection{Angular correlation functions}
We also try using three different functions for $\chiIJhd$ for the time reallocation as maximising equation~\ref{eqn:snrPIPJPg} may not totally capture a desired observing strategy for PTAs. 


\subsubsection{Hellings-Downs}
The Hellings-Downs function, Eqn.~\ref{eqn:hd} is how we expect the gravitational wave signal to appear in pulsar pairs. 
It is shown as the blue line inf Fig.~\ref{fig:correlations}.
It is important for a PTA to have good coverage of a range of angular separations between pulsar pairs so that a gravitational-wave detection can be confirmed. 
However, if we use the Hellings-Downs function in our SNR calculation, we find that we favour pulsars in the galactic centre where the angular separations are small and the value of Eqn.~\ref{eqn:hd} is largest. 

\subsubsection{Dipole $-$ Hellings-Downs}
As we want to keep pulsar pairs with large angular separations we also try maximising a modified Eqn.~\ref{fig:correlations} with a different $\chiIJhd$. 
One option is to use the difference between Hellings-Downs function and a dipole, 
%Eqn.~\ref{eqn:hd} has highest values for pulsars that have small angular separations.
%But we want good sky coverage to confirm a Hellings-Downs correlation for detection. 
%To prevent the time-shuffle algorithm favouring pulsars in the galactic centre, we replace the Hellings and Downs correlation $\chiIJhd$ with a modified angular correlation function $\chi_{IJ}^{\rm mod}$ which corresponds to the difference between a dipole $\chiIJdip$ and the HD correlation
\begin{equation}
\chi_{IJ}^{\rm mod} = \chiIJdip - 2 \chiIJhd
\label{eqn:chimod}
\end{equation}
where $\chiIJdip$ is given by~\cite{AnholmEtAl:2009}
\begin{equation}
\chiIJdip = -\frac{3}{2}
            \left[ \cos \zeta - \frac{4}{3} - 4 \tan^2 \frac{\zeta}{2} 
                   \ln \left( \sin \frac{\zeta}{2} \right) \right]\,.
\label{eqn:dip}
\end{equation}
The resulting function now favours large angular separations between pulsars and is shown as the green line in Fig.~\ref{fig:correlations}. 

\subsubsection{Equal value}
A third option is to have no preference for any angular separation and set $\chiIJhd$ to a constant value for all pulsar pairs. 



\begin{figure}
\includegraphics[width=0.5\textwidth]{../snr/correlationFunctions.png}
\caption{\label{fig:correlations}
Correlation functions. 
The blue line shows the Hellings and Downs curve, Eqn.~\ref{eqn:hd}. 
The orange line shows a dipole correlation, Eqn.~\ref{eqn:dip} and the green line shows a correlation function which no longer favours the galactic centre pulsars (Eqn.~\ref{eqn:chimod}).
The grey markers at the bottom of the plot indicate the angular coverate of the pulsars pairs being timed by MeerTime.
}
\end{figure}

\subsection{Results}

Results of our time reallocation method are shown in Fig.~\ref{fig:snrCompare}.
The original observation stragegy is shown by the blue-solid line.
The results of the three different angular correlation functions are shown.
The orange-dotted line shows the results of using the Hellings-Downs function. 
The green-dasshed line shows the result of using the dipole minus Helllings-Downs function. 
The red-dot-dashed line shows the results of using all equal values for angular correlation. 
(The different functions are only used for the time reallocatoin algoritm, the plot shows the actual SNR value using the Hellings-Downs function as in Eqn.~\ref{eqn:snrPIPJPg}).

Using the Hellings-Downs function is the best performing, but comes at the cost of losing pulsar pairs at high angular separation.
The equal-values is the next best and provides and SNR improvement of $\approx 1$ after $10\,{\rm yrs}$.



\begin{figure}
\includegraphics[width=0.5\textwidth]{../local/rjNoiseSoFar/compareSNRVTime.png}
\caption{\label{fig:snrCompare}
SNR over PTA lifespan ($T$) assuming four different observation time allocations. 
Blue-solid shows the original MeerTime observing strategy.
Three time realocation results are shown, which each assume a different function for $\chiIJhd$. 
`hd' is Hellings-Downs, `dphdDiff' is the difference between dipole and Hellings-Downs, and `equal' is equal value for all pulsar pairs. 
}
\end{figure}

In this work we take a greedy algorithm approach, which will find the optimal step at each stage, however it is not guarenteed to find the global optimal solution. 

\section{Some notes / thoughts}
This method is definitely not finding an optimal solution, I think it finds a local maximum in SNR, but we start from a reasonable original strategy and try to improve on that, so it may not be so bad. 
At the moment, I am only changing the observation time allocation for each pulsar, so the only parameter we are changing is $\sigma_I$ which depends on $T_{{\rm obs},I}$.
We could also think about changing the cadence. 

In terms of pulsar timing data for the future, people will also be thinking about the cost of losing data for a particular pulsar.  
For example is it beneficial enough to remove a pulsar from the observation schedule and lose the opportunity for decadels-long observations and study of individual objects for the benefit of gravitational-wave detection? 



\bibliographystyle{abbrv}
\bibliography{ptabib} % if your bibtex file is called example.bib


\end{document}
